%% ---------------------------------------------------------------------------------------
\section{Minimising the size of the network}
\label{sec:min}

``Size'' can be minimised at four levels; we tried all four and report parameters (or non-zero
weights) of everything below the trunk.

\paragraph{1. Design choices that shrink the tree at equal (100\%) training accuracy.}
(a) \emph{Node width} has an interior optimum (Table~\ref{tab:width}): linear or width-2 children
make no progress and fall back to isolation nodes that need one 129-parameter cap per memorised
image (0.84--1.03M parameters), very wide children are expensive per node (259k at $h=128$), and
$h=8$--32 gives the smallest trees (78--80k). (b) \emph{Sharing the root} across classes: the corrector
tree needs half the parameters of the one-vs-rest forest (Table~\ref{tab:archs}). (c) \emph{Letting
each node choose its trunk} gives the largest reduction we found: 144k$\to$44k parameters
(83$\to$35 nodes) on CIFAR-10 at higher test accuracy (Table~\ref{tab:main}).

\paragraph{2. Decision-preserving pruning.} For every node we search the highest sparsity
$s$ among 0.5, 0.8, 0.9, 0.95, 0.98 and 0.99 at which magnitude pruning~\cite{lottery} followed by a short masked
fine-tune leaves the node's decision on \emph{every one of its training samples} unchanged.
Because a node's decisions define its children's training sets (Table~II), this preserves the whole
tree's 100\% training accuracy by construction; only the connection weights are recomputed with
\eqref{eq:closed}. Table~\ref{tab:prune} shows 20--32\% fewer non-zero weights at unchanged training
accuracy and test accuracy within $\pm0.1$ points. The prunability is very uneven: nodes that serve a
few hundred samples prune to 80--95\% sparsity, while the root and the first-level detectors (which
must keep 45--50k decisions exactly) cannot be pruned even to 50\%. Pruning also performs the note's
\emph{feature selection} (open question 5): the sparsest nodes read only 34--40 of the 128 PCA features.

\begin{table}[h]
\centering\small
\caption{Decision-preserving pruning: every node keeps every training decision, so the tree stays at
100\% training accuracy.}
\label{tab:prune}
\resizebox{\linewidth}{!}{\begin{tabular}{@{}lrrrrrr@{}}\toprule
tree & non-zero weights & after pruning & reduction & train acc. & test acc. before & after\\\midrule
CIFAR-10, ResNet, corrector $h$=8 & 77\,612 & 58\,422 & 25\% & 100.0\% & 69.78\% & 69.75\%\\
CIFAR-10, CNN, corrector $h$=8 & 144\,044 & 112\,116 & 22\% & 100.0\% & 81.81\% & 81.92\%\\
CIFAR-10, ResNet, corrector $h$=32 & 79\,643 & 55\,289 & 31\% & 100.0\% & 70.75\% & 70.79\%\\
Cats vs Dogs, CNN, binary $h$=8 & 17\,430 & 11\,808 & 32\% & 100.0\% & 92.06\% & 92.02\%\\
Cats vs Dogs, ResNet, binary $h$=8 & 5\,205 & 4\,173 & 20\% & 100.0\% & 93.24\% & 93.24\%\\
\bottomrule\end{tabular}
}
\end{table}

\paragraph{3. The smallest interpolating network sits at the double-descent peak.}
If ``minimise the size'' is taken literally -- the fewest parameters subject to 100\% training
accuracy -- the double-descent sweep answers it (Table~\ref{tab:smallest}): the smallest interpolating
model is a \emph{single} MLP of exactly the threshold width (2.4k parameters with clean labels,
4.8k with noise), smaller than any tree, and it is the \emph{worst} of all interpolating MLPs on test data
(37.7\% and 53.8\% error). Trees are 3--5$\times$ larger but 3--6 points better, and the best
interpolating model is $64\times$ larger than the smallest one. Minimum size and good generalisation are
opposite goals once 100\% training accuracy is imposed.

\begin{table}[h]
\centering\small
\caption{Interpolating models in the double-descent sweep (CIFAR-10, $n=4\,000$). The best test
error at any width is 30.2\% (clean; an interpolating MLP) and 33.2\% (noisy; a non-interpolating MLP, $H=8$).}
\label{tab:smallest}
\begin{tabular}{@{}llrr@{}}\toprule
labels & interpolating model & parameters & test error\\\midrule
clean & smallest single MLP ($H$=32) & 2\,410 & 37.7\%\\
 & smallest tree ($H$=8; 23 nodes) & 12\,248 & 34.9\%\\
 & best interpolating model (MLP, $H$=2048) & 153\,610 & 30.2\%\\
\midrule
20\% noisy & smallest single MLP ($H$=64) & 4\,810 & 53.8\%\\
 & smallest tree ($H$=8; 31 nodes) & 16\,480 & 48.2\%\\
 & best interpolating model (MLP, $H$=4096) & 307\,210 & 39.0\%\\
\bottomrule\end{tabular}

\end{table}

\paragraph{4. Cost-complexity pruning of subtrees.} Dropping the 100\% requirement, we repeatedly
remove the subtree that fixes the fewest training samples per parameter (as in CART pruning).
Fig.~\ref{fig:prune} traces the resulting path from the full tree to the bare root: test accuracy
changes almost linearly with training accuracy, losing 0.3--0.5 points of test accuracy for every
point of training accuracy the subtrees add. On these datasets the best tree for \emph{test}
accuracy is the root alone, and on Cats vs Dogs (CNN) a 1\,041-parameter root achieves 94.1\% test
accuracy, 2 points better than the 17-node interpolating tree.

\begin{figure}[h]
\centering
\includegraphics[width=0.55\linewidth]{prune_curve.pdf}
\caption{Cost-complexity pruning, read from right (full tree, 100\% training accuracy) to left (root
only). Each point removes one more subtree.}
\label{fig:prune}
\end{figure}
